\documentclass[10pt,a4paper]{article}
\usepackage[margin=0.6in]{geometry}
\usepackage{amsmath,amssymb,amsfonts,stmaryrd}
\usepackage{xcolor}
\usepackage{enumitem}
\usepackage{tcolorbox}
\usepackage{multicol}
\usepackage{listings}
\usepackage{hyperref}

\definecolor{isabelleblue}{RGB}{0, 70, 140}
\definecolor{isabellekw}{RGB}{120, 0, 120}
\definecolor{boxbg}{RGB}{248, 249, 250}

\tcbset{
  colback=boxbg,
  colframe=isabelleblue,
  fonttitle=\bfseries\large,
  boxrule=0.8pt,
  arc=3pt,
  left=8pt,
  right=8pt,
  top=6pt,
  bottom=6pt
}

\title{\textbf{\huge Proofs by Substitution}\\[0.3em]
\large Lecture Notes Transcription}
\author{\textbf{Course:} Interactive Theorem Proving (7CCSACTP)}
\date{}

\begin{document}

\maketitle
\vspace{-1.5em}

\begin{tcolorbox}[title=Context \& Concept of Substitution]
\begin{multicols}{2}
\textbf{$\rightarrow$ Previous Videos:}
\begin{itemize}[leftmargin=1.2em, itemsep=2pt]
    \item Introduced Isabelle interface
    \item How to write terms / theorems in Isabelle
    \item Two basic proof / reasoning approaches:
    \begin{itemize}[leftmargin=1em, label=$\bullet$]
        \item Forward reasoning (FW)
        \item Backward reasoning (BW)
    \end{itemize}
\end{itemize}

\columnbreak

\textbf{$\rightarrow$ \underline{This Video}: Proofs by Substitution}
\begin{itemize}[leftmargin=1.2em, itemsep=2pt]
    \item Suppose we have a theorem:
    \begin{center}
    \texttt{lem1:} $\{?x + ?y = ?y + ?x\} \longrightarrow$ \textit{simplification rule}
    \end{center}
    \item And we have a goal: \quad \texttt{goal :: 1 + 3 > 3 + 1}
    \item Substituting \texttt{lem1} into the goal gives a new goal:
    \begin{center}
    \texttt{goal :: 3 + 1 > 3 + 1}
    \end{center}
    \item \textbf{When does substitution make sense?}
    \begin{itemize}[leftmargin=1em, label=$\bullet$]
        \item If R.H.S. of the simp rule is simpler than L.H.S., or if we have a proof on R.H.S.
    \end{itemize}
\end{itemize}
\end{multicols}
\end{tcolorbox}

\vspace{0.3em}

\begin{tcolorbox}[title=Worked Example: Distributivity on Nats]
\textbf{Goal:} $\vdash (x + y) \cdot (a + b) = x \cdot a + x \cdot b + y \cdot a + y \cdot b \qquad (\text{Note: all variables are } \text{nat})$

\begin{enumerate}[leftmargin=1.5em, itemsep=4pt]
    \item $(x + y) \cdot (a + b) = (x + y) \cdot a + (x + y) \cdot b$
    \begin{itemize}[leftmargin=1em, label=$\rightarrow$]
        \item \textit{Using rule:} $?x \cdot (?a + ?b) = ?x \cdot ?a + ?x \cdot ?b$ \quad (\texttt{lem1})
    \end{itemize}
    
    \item $\underline{(x + y) \cdot a} + (x + y) \cdot b = x \cdot a + y \cdot a + (x + y) \cdot b$
    \begin{itemize}[leftmargin=1em, label=$\rightarrow$]
        \item \textit{Using rule:} $(?a + ?b) \cdot ?x = ?a \cdot ?x + ?b \cdot ?x$ \quad (\texttt{lem2})
    \end{itemize}
    
    \item $\dots = x \cdot a + y \cdot a + x \cdot b + y \cdot b \quad \checkmark$
    \begin{itemize}[leftmargin=1em, label=$\rightarrow$]
        \item \textit{Using distributive and associative lemmas.}
    \end{itemize}
\end{enumerate}
\end{tcolorbox}

\vspace{0.3em}

\begin{tcolorbox}[title=Isabelle Practical Notes]
\begin{itemize}[leftmargin=1.2em, itemsep=3pt]
    \item \textbf{Show proof steps:} How to perform the proof step-by-step using rules.
    \item \textbf{Library Search:} Find lemmas that already exist in Isabelle's library.
    \item \textbf{Reflexivity:} A goal of the form $x = x$ can be proved using \texttt{by (rule refl)} or \texttt{by simp}.
    \item \textbf{Associativity convention:} $a + b + c$ is parsed as $a + (b + c)$.
    \item \textbf{Symmetric attribute:} If \texttt{lem1: x = y}, then \texttt{lem1[symmetric]} gives \texttt{y = x}.
    \item \textbf{Closing proofs:} If the goal is $P \vdash P$, close it by \texttt{"."} (or \texttt{by assumption}).
    \item \textbf{Calculational Reasoning (Isar):}
    \begin{center}
    \begin{tabular}{lll}
    \texttt{have} & \texttt{"x = y"} & \texttt{...} \\
    \texttt{also have} & \texttt{"... = z"} & \texttt{...} \\
    \texttt{also have} & \texttt{"... = a"} & \texttt{...} \\
    \texttt{also have} & \texttt{"... = b"} & \texttt{...} \\
    \texttt{finally show} & \texttt{"x = b"} & \texttt{.}
    \end{tabular}
    \end{center}
\end{itemize}
\end{tcolorbox}

\end{document}
